\section{核岭回归与核方法的代价}
\label{sec:13-Machine-Learning/09-Kernel-Methods-and-Gaussian-Processes/04}

同一份数据，同一台 64 GB 内存的机器。线性 Ridge 两分钟跑完，换成 RBF 核之后，脚本在构造核矩阵那一行被系统直接杀掉。样本数是十二万，特征只有四十维。奇怪的地方在于：核方法号称把维度的账单抹掉了，这里维度本来就不高，它反而先撑不住。

这一节先兑现上一节的许诺——把表示定理用于平方损失，推出核岭回归的闭式解；再把它与\hyperref[sec:13-Machine-Learning/02-Linear-Models/02]{``正则化从 Ridge 到 Lasso''}中的原始形式对照，看清同一件事的两种视角；最后把那台机器上发生的事算成三笔账。结论会很直白：核化不是免费升维，而是用样本空间中的二次存储与三次求解，换特征空间的表达力。

\subsection{平方损失下解可以写出来}

给定训练数据 \((x_1,y_1),\dots,(x_n,y_n)\)、核 \(k\) 与正则参数 \(\lambda>0\)。核岭回归求解

\[
\min_{f\in\mathcal{H}_k}\
\sum_{i=1}^{n}\bigl(y_i-f(x_i)\bigr)^{2}
+\lambda n\,\norm{f}_{\mathcal{H}_k}^{2}.
\]

写成 \(\lambda n\) 是为了让闭式解干净：损失是求和而非平均，正则强度随 \(n\) 缩放。由表示定理，极小解必为 \(f=\sum_{j=1}^{n}\alpha_j k(x_j,\cdot)\)。记 Gram 矩阵 \(K=\bigl[k(x_i,x_j)\bigr]\in\R^{n\times n}\)，则

\[
f(x_i)=\sum_{j=1}^{n}K_{ij}\alpha_j=(K\alpha)_i,
\qquad
\norm{f}_{\mathcal{H}_k}^{2}=\alpha^{\trans}K\alpha,
\]

问题化为 \(n\) 维二次问题

\[
\min_{\alpha\in\R^{n}}\
\norm{y-K\alpha}^{2}+\lambda n\,\alpha^{\trans}K\alpha.
\]

这是 \(\alpha\) 的凸二次函数：第一项的 Hessian 是 \(2K^{2}\succeq0\)，第二项的是 \(2\lambda nK\succeq0\)，凸性来自 \(K\) 半正定，也就是上一节反复强调的那个条件。令梯度为零：

\[
-2K(y-K\alpha)+2\lambda nK\alpha=0
\iff
K\bigl[(K+\lambda nI)\,\alpha-y\bigr]=0.
\]

解括号内的方程即为一个充分取法。由于 \(K\succeq0\) 且 \(\lambda n>0\)，矩阵 \(K+\lambda nI\) 对称正定、必然可逆，于是

\[
\widehat\alpha=(K+\lambda nI)^{-1}y,
\qquad
\widehat f(x)=\sum_{i=1}^{n}\widehat\alpha_i\,k(x_i,x).
\]

\(K\) 奇异时梯度方程还有别的解，但它们只在 \(K\) 的零空间方向上相差一个分量，给出的函数相同，上面的闭式解总是合法的。训练归结为一次 \(n\times n\) 线性方程组求解，预测是 \(n\) 个核值的加权和。

\subsection{原始与对偶是同一个解的两种坐标}

取线性核 \(k(x,z)=\ip{x}{z}\)，此时 \(K=XX^{\trans}\)，闭式解成为 \(\widehat\alpha=(XX^{\trans}+\lambda nI)^{-1}y\)，预测

\[
\widehat f(x)=\sum_{i=1}^{n}\widehat\alpha_i\,\ip{x_i}{x}
=\ip{X^{\trans}\widehat\alpha}{x}.
\]

令 \(\widehat w=X^{\trans}\widehat\alpha\)，与 Ridge 的原始解 \(\widehat w=(X^{\trans}X+\lambda nI)^{-1}X^{\trans}y\)（正则参数按同一尺度写出）对比，两者给出完全相同的预测——恒等式

\[
X^{\trans}(XX^{\trans}+\lambda nI)^{-1}
=(X^{\trans}X+\lambda nI)^{-1}X^{\trans}
\]

两端右乘 \(y\) 即见。差别只在于解哪个方程组，而这一点由矩阵的形状决定。

\begin{center}
\begin{tikzpicture}[font=\small]

\fill[black!8] (0,0) rectangle (1.1,3.4);
\draw (0,0) rectangle (1.1,3.4);
\node[anchor=north,font=\footnotesize,align=center] at (0.55,-0.15)
  {\(X\)\\\(n\times p\)};

\fill[black!8] (3.0,2.3) rectangle (4.1,3.4);
\draw (3.0,2.3) rectangle (4.1,3.4);
\node[anchor=north,font=\footnotesize,align=center] at (3.55,2.1)
  {\(X^{\trans}X+\lambda nI\)\\\(p\times p\)};
\node[anchor=west,font=\footnotesize] at (4.35,2.85) {原始形式};

\fill[black!8] (7.2,0) rectangle (10.6,3.4);
\draw (7.2,0) rectangle (10.6,3.4);
\node[anchor=north,font=\footnotesize,align=center] at (8.9,-0.15)
  {\(K+\lambda nI\)\\\(n\times n\)};
\node[anchor=south,font=\footnotesize] at (8.9,3.55) {核（对偶）形式};

\draw[->,>=stealth,black!55] (1.35,2.85) -- (2.75,2.85);
\draw[->,>=stealth,black!55] (1.35,1.0) -- (6.95,1.0);
\node[font=\footnotesize,black!60,fill=white,inner sep=2.5pt] at (4.15,1.0)
  {与 \(p\) 无关，只看 \(n\)};

\end{tikzpicture}
\end{center}

\emph{同一批数据，两种形状的方程组；哪一边划算，取决于 \(n\) 与 \(p\) 谁更大。}

原始形式在特征权重里思考：未知量 \(w\in\R^{p}\)，求解 \(p\times p\) 系统，适合特征不多而样本很多的情形。对偶形式在样本相似度里思考：未知量 \(\alpha\in\R^{n}\)，求解 \(n\times n\) 系统，特征维数 \(p\) 从账单中彻底消失。多项式核、RBF 核对应的 \(p\) 巨大甚至无穷，原始形式根本写不出来，对偶形式却只用 \(n\times n\) 的 \(K\)，这正是``不构造 \(\phi\)''的兑现。图里右边那块方阵在 \(n\) 大时会长到吓人，开头那台机器倒在的就是它。

\subsection{两个旋钮控制的不是同一件事}

核岭回归有两个主要超参数，作用机制完全不同。

\(\lambda\) 控制范数预算。正则项 \(\lambda n\,\alpha^{\trans}K\alpha\) 直接压在 \(\norm{f}_{\mathcal{H}_k}\) 上：\(\lambda\to0\) 时模型几乎只追求拟合，若 \(K\) 正定则退化为插值，训练误差为零；\(\lambda\to\infty\) 时 \(\widehat\alpha\to0\)，模型塌成常数。中间地带是偏差与方差的权衡，与\hyperref[sec:13-Machine-Learning/02-Linear-Models/02]{``正则化从 Ridge 到 Lasso''}中 \(\lambda\) 的角色一致，只是这里收缩的是 RKHS 几何下的不光滑方向——对应 \(K\) 的小特征值方向——而非坐标方向的系数。

RBF 核的 \(\gamma\) 控制的是长度尺度，也就是核所声明的``什么算平滑''。\(\gamma\) 大时 \(k(x,z)\) 只在紧邻处非零，\(K\) 接近单位阵，几乎每个样本自立锚点，函数可以紧贴数据剧烈起伏；\(\gamma\) 小时远处样本也彼此相似，\(K\) 接近全一矩阵，允许的函数接近常数。两个旋钮联合决定有效容量：大 \(\gamma\) 配大 \(\lambda\) 是``精细但被强压''的模型，小 \(\gamma\) 配小 \(\lambda\) 是``本来就粗''的模型。验证误差的表面因此是二维的，只沿一条轴调参会漏掉谷底。

选择它们的纪律与\hyperref[sec:11-Mathematical-Statistics/06-Resampling-and-Model-Selection/03]{``交叉验证评估的是完整训练流程''}强调的完全相同：预处理、核的选择与超参数调节都必须包在每个训练折内部，用嵌套验证或从未参与选择的测试集评估整个流程。核化没有改变评估纪律，只改变了被评估的对象。

\subsection{Gram 矩阵的三笔账}

核方法把特征维数从账单上抹掉，代价是把样本量写了上去：

\begin{itemize}
\tightlist
\item
  存储 \(K\)：\(O(n^{2})\)。\(n=10^{5}\) 时 \(K\) 有 \(10^{10}\) 个双精度数，约 80 GB，单台机器已放不下；
\item
  求解 \((K+\lambda nI)\alpha=y\)：稠密 Cholesky 分解 \(O(n^{3})\)。\(n=10^{5}\) 约需 \(10^{15}\) 次浮点运算；
\item
  预测：每个新样本要算 \(n\) 个核值，\(O(nd)\) 起步，且全部训练数据必须随模型保存。
\end{itemize}

开头那台机器上，十二万样本对应的 \(K\) 约 115 GB，第一笔账就付不起，与特征只有四十维毫无关系。这三笔账也解释了朴素核方法为何在大数据时代让位于参数模型：\(O(n^{3})\) 的求解意味着样本量翻倍，训练时间变成八倍。

两条常用出路都在``核给的表达力''与``线性方法的规模''之间找折中。Nyström 近似抽样 \(m\ll n\) 个诱导点，只用 \(K\) 的对应列做低秩近似 \(K\approx CW^{-1}C^{\trans}\)，把存储与求解降到 \(O(nm)\) 与 \(O(m^{3})\) 量级；随机特征则对平移不变核显式构造低维特征 \(z(x)\in\R^{m}\)，使 \(\ip{z(x)}{z(x')}\approx k(x,x')\)，然后回到原始形式训练线性模型——核被重新显式化，训练成本回到线性模型的水平。两条路的近似质量都由 \(m\) 控制，\(m\) 取小了，被丢掉的正是核谱里衰减慢的那些方向。

\subsection{为什么还要付这笔账}

如果只为拟合训练数据，核方法并非必需；为它定价的是泛化。\hyperref[ch:118]{``统计学习理论：训练误差对真风险有多乐观''}一章的核心问题是训练误差比真风险乐观多少，答案取决于所用函数类的有效容量。对核方法，容量恰由范数球

\[
\bigl\{\,f\in\mathcal{H}_k:\ \norm{f}_{\mathcal{H}_k}\le B\,\bigr\}
\]

刻画：这类函数集的复杂度量（如 Rademacher 复杂度）以 \(B\) 与核的谱衰减速度为界，\(\lambda\) 通过限制 \(B\) 直接收缩训练误差与真风险的差距。正则项因此不是工程技巧，\(\norm{f}_{\mathcal{H}_k}\) 正是控制有效复杂度的那个范数——这与 SVM 中``最大化间隔等价于最小化 \(\norm{w}^{2}\)''是同一枚硬币的两面。

到这里，本单元开头那个观察已经走得相当远：数据只以内积出现，于是有核技巧；核技巧需要一个家，于是有 RKHS 与表示定理；表示定理落地为闭式解，标价 \(O(n^{2})\) 与 \(O(n^{3})\)。它走不到的地方也标清楚了——大样本的账单，以及``核由人来选''这个没有被任何定理消除的自由度。下一节换一个视角看同一个 \(K\)：把它读成函数值的联合协方差，那个自由度就成了先验的选择，预测也从一个数变成一个分布。
